\documentclass{ximera}

\title{Infinite Limits and Indeterminate Forms Activity Facilitator Guide}
\author{MATH 425: Calculus I}

\begin{document}
\begin{abstract}
    Working with the peers in your group, solve the following problems. Make sure to show and justify all your work. Make sure everyone in the group understands the solution and participates. Be prepared to report your answers to the whole class. 
\end{abstract}
\maketitle

\textbf{Student Learning Objectives}
Students will be able to:
\begin{itemize}
    \item Evaluate limits of $\frac{c}{0}$ form where $c\neq0$ using one sided limits.
    \item Reason using definitions of limits, continuity and existance of limits.
    \item Evalute indeterminate limits using multiplication by the conjugate.
    \item Evaluate indeterminate limits using methods of combining fractions.
    \item Identify locations where limits of piecewise functions do not exist.
\end{itemize}
\textbf{Prerequisite Knowledge:} This activity expects students to be familiar with basic limits, one-sided limits and continuity. Familiarity with processes of using one-sided limits and multiplying by a conjugate are also expected.  

\begin{exercise}
    Determine whether the following statements are true or false.  Explain your reasoning.
    \begin{enumerate}
        \item If we state that $\lim_{x\rightarrow 5}f(x)=\infty$, then we are implicitly saying that the limit exists.\\
        \textcolor{blue}{False: A limit existing implies that the limit has a finite value.  Saying a limit equals $\infty$ is a specific way that a limit can fail to exist.}\\
        \textcolor{orange}{This is something we try to reinforce with students in class, but the notation is misleading.  Helpful verbage for talking about a limit of infinity is that it is a ``trend'' rather than an exact value.  This is a good time to start making sure students understand that $\infty$ is NOT a number and we cannot do algebra with it. }
        \item If we state that $\lim_{x\rightarrow 5}f(x)=5$, then we are implicitly saying that the limit exists.\\
        \textcolor{blue}{True: This limit is a finite, numerical value.}
        \item If $\lim_{x\rightarrow 1^-}f(x)=\infty$, then $\lim_{x\rightarrow 1^+}f(x)=-\infty$.\\
        \textcolor{blue}{False: This does not need to be true.  For example $f(x)=\frac{1}{(x-1)^2}$ has $\lim_{x\rightarrow 1^-}f(x)=\infty$ and $\lim_{x\rightarrow 1^+}f(x)=\infty$.}
        %\item If $\lim_{x\rightarrow 5}f(x)=\infty$, then $f(x)$ has a vertical asymptote at $x=5$.
        \item If $\lim_{x\rightarrow 5}f(x)=\infty$, then $f(x)$ is continuous at $x=5$.\\
        \textcolor{blue}{False: Continuity requires the limit to exist at $x=5$, and this limit does not exist.}\\
         Problems adapted from Hartman, et al., \textit{APEX Calculus}.
    \end{enumerate} 
\end{exercise}

\begin{exercise}
    Evaluate the limits, including classifying infinite limits (one-sided and two-sided, where appropriate).
\begin{enumerate}
    \item $\displaystyle \lim_{x\rightarrow -2}\frac{\sqrt{x^2+21}-5}{x+2}$ \\
    \textcolor{blue}{ Gives a $\frac{0}{0}$ form.\\
    $\displaystyle=\lim_{x\rightarrow -2} \frac{\sqrt{x^2+21}-5}{x+2}\bigg(\frac{\sqrt{x^2+21}+5}{\sqrt{x^2+21}+5}\bigg)$\\
    $\displaystyle=\lim_{x\rightarrow-2} \frac{x^2+21-25}{(x+2)(\sqrt{x^2+21}+5)}$\\
    $\displaystyle=\lim_{x\rightarrow -2}\frac{(x-2)(x+2)}{(x+2)(\sqrt{x^2+21}+5)}$\\
    $\displaystyle=\lim_{x\rightarrow -2}\frac{x-2}{\sqrt{x^2+21}+5}=\frac{-4}{10}=-\frac25$}
    
    \item $\displaystyle \lim_{x\rightarrow -1}\frac{x-2}{x^2+2x+1}$ \\
    \textcolor{blue}{ Gives a $\frac{c}{0}$ form\\
    numerator: $-1-2=-3$\\
    denominator: if $x>-1$, $x^2+2x+1\rightarrow 0$ and $(x+1)^2>0$\\
    if $x<-1$. $x^2+2x+1\rightarrow 0$ and $(x+1)^2>0$\\
    Either way, the fraction is negative\\
    $\displaystyle\lim_{x\rightarrow -1}\frac{x-2}{x^2+2x+1}=-\infty$}
    
    
    \item $\displaystyle\lim_{x\rightarrow 0}\frac{\sqrt{x^2+4}-2}{x^2}$ \\
    \textcolor{blue}{ Gives $\frac00$ form\\
    $\displaystyle\lim_{x\rightarrow 0}\frac{\sqrt{x^2+4}-2}{x^2}\bigg(\frac{\sqrt{x^2+4}+2}{\sqrt{x^2+4}+2}\bigg)$\\
    $\displaystyle =\lim_{x\rightarrow 0}\frac{x^2+4-4}{x^2(\sqrt{x^2+4}+2)}$\\
    $=\displaystyle =\lim_{x\rightarrow 0}\frac{1}{\sqrt{x^2+4}+2}=\frac{1}{4}$}
    

    \item $\displaystyle \lim_{x\rightarrow 0}\frac{3x-2}{x^2-x}$\\
     \textcolor{blue}{ Gives a $\frac{c}{0}$ form\\
   numerator: $3(0)-2=-2$\\
    denominator: As $x\rightarrow 0$, $x^2\leq x$.  If $x<0$ and $x\rightarrow 0$, $x^2-x>0$ (fraction is negative).\\
    If $x>0$ and $x\rightarrow 0$, $x^2-x<0$ (fraction is positive).\\
    Thus $\displaystyle \lim_{x\rightarrow 0^-} \frac{3x-2}{x^2-x}=-\infty$ and $\displaystyle \lim_{x\rightarrow 0^+} \frac{3x-2}{x^2-x}=\infty$.\\
    The one-sided limits do not match, so the limit DNE.}\\
    \textcolor{orange}{There is some sophisticated thinking involved in figuring out whether the denominator is positive or negative here.  Students can either use the fact that near 0, $x^2<x$, or they can factor $x^2-x=x(x-1)$ and determine that if $x$ is near 0, $x-1$ has to be negative.}

    \item $\displaystyle \lim_{x\rightarrow 3} \frac{1-x}{|x-3|}$\\
    \textcolor{blue}{ numerator: $1-3=-2$\\
    denominator: Regardless of whether $x$ is coming from the left or right, $|x-3|>0$ as $x\rightarrow 3$.  Thus the fraction is negative in both cases, or $\displaystyle \lim_{x\rightarrow 3^-}\frac{1-x}{|x-3|}=\lim_{x\rightarrow 3^+}\frac{1-x}{|x-3|}=-\infty=\lim_{x\rightarrow 3}\frac{1-x}{|x-3|}$.}

    \item $\displaystyle \lim_{x\rightarrow 0} \bigg( \frac{1}{x}-\frac{1}{x(x-1)}\bigg)$\\
    \textcolor{blue}{ $=\displaystyle\lim_{x\rightarrow 0} \frac{(x-1)}{x(x-1)}-\frac{1}{x(x-1)}=\lim_{x\rightarrow 0} \frac{x-2}{x(x-1)}$ gives a $\frac{c}{0}$ form.\\
    Numerator: $0-2=-2$\\
    Denominator: if $x<0$, then $x(x-1)\rightarrow 0$ and is positive (fraction is negative).\\
    If $x>0$, then $x(x-1)\rightarrow 0$ and is negative (fraction is positive)\\
    Limit DNE}

     \item $\displaystyle \lim_{x\rightarrow 0} \bigg( \frac{1}{x}+\frac{1}{x(x-1)}\bigg)$\\
     \textcolor{blue}{ $=\displaystyle\lim_{x\rightarrow 0} \frac{(x-1)}{x(x-1)}+\frac{1}{x(x-1)}=\lim_{x\rightarrow 0} \frac{x}{x(x-1)}=\lim_{x\rightarrow0} \frac{1}{x-1}=\frac{1}{-1}=-1$}


  \end{enumerate}
\end{exercise} 

\begin{exercise}
    Consider the function \[ f(x)=
    \begin{cases}
    3x^2-1 & \text{if }x<-1\\
    -2x & \text{if } -1<x<2\\
    \frac{2}{x-3} &\text{if } 2<x
    \end{cases}\]
    Identify all $a$ where $\displaystyle \lim_{x\rightarrow a} f(x)$ does not exist.  Show reasoning to support your answer.\\
    \textcolor{blue}{ This piecewise function changes at $x=-1$ and $x=2$, so the limit may fail to exist at these points. Also, $x=3$ causes division by 0.  Each of the rules is defined everywhere else.\\
    $x=-1$:\\
    $\displaystyle \lim_{x\rightarrow -1^-} f(x)=3(-1)^2-1=2$\\
    $\displaystyle \lim_{x\rightarrow -1^+} f(x)=-2(-1)=2$\\
    So $\displaystyle \lim_{x\rightarrow-1}f(x)=2$\\
    $x=2$:\\
    $\displaystyle \lim_{x\rightarrow 2^-} f(x)=-2(2)=-4$\\
    $\displaystyle \lim_{x\rightarrow -2^+} f(x)=\frac{2}{2-3}=-2$\\
    So $\displaystyle \lim_{x\rightarrow 2}f(x)$ DNE\\
    $x=3$:\\
    numerator: 2\\
    denominator: If $x<3$, $x-3<0$ so $\displaystyle \lim_{x\rightarrow 3^-} f(x)=-\infty$\\
    If $x>3$, $x-3>0$ so $\displaystyle \lim_{x\rightarrow 3^+}f(x)=\infty$\\
    $\displaystyle \lim_{x\rightarrow 3} f(x)$ DNE\\
    \\
    So $\displaystyle \lim_{x\rightarrow a} f(x)$ DNE when $a=2,3$
    }\\
    \textcolor{orange}{Please check in with students to make sure that they are not just ``subbing in'' to both sides to get the answers at the break points.  We really want to reinforce the idea of using one-sided limits here.  Also, some students may need to be gently reminded that there can be other ``issue'' points besides just at the endpoints of the domain sections.}
\end{exercise}


\begin{exercise}

  Challenge problem: Evaluate $\displaystyle \lim_{x\rightarrow 2} \frac{\sqrt{6-x}-2}{\sqrt{3-x}-1} $\\
  \textcolor{blue}{ $=\displaystyle \lim_{x\rightarrow 2} \frac{\sqrt{6-x}-2}{\sqrt{3-x}-1}\bigg(\frac{\sqrt{6-x}+2}{\sqrt{6-x}+2}\bigg)\bigg(\frac{\sqrt{3-x}+1}{\sqrt{3-x}+1}\bigg)$\\
$=\displaystyle \lim_{x\rightarrow 2}\frac{(6-x-4)(\sqrt{3-x}+1)}{(3-x-1)(\sqrt{6-x}+2)}$\\
$=\displaystyle \lim_{x\rightarrow 2}\frac{(2-x)(\sqrt{3-x}+1)}{(2-x)(\sqrt{6-x}+2)}$\\
$=\displaystyle \lim_{x\rightarrow 2}\frac{\sqrt{3-x}+1}{\sqrt{6-x}+2}$\\
$\displaystyle =\frac{\sqrt{3-2}+1}{\sqrt{6-2}+2}=\frac{2}{4}=\frac12$}\\
\textcolor{orange}{This is definitely harder than would appear on an exam or quiz, but really reinforces the conjugate for students.}
     

\end{exercise}



\end{document}